\section{Gravedad unimodular}
\label{sec:ug}

\subsection{La idea}

El punto de partida es el problema de la constante cosmológica que se mencionó en la sección~\ref{sec:rg}. En RG, la energía del vacío de los campos de materia gravita exactamente igual que una constante cosmológica desnuda, porque entra en el lado derecho de \eqref{eq:Einstein} como un término proporcional a $g_{ab}$. La gravedad unimodular propone cambiar la teoría de tal modo que ese término deje de gravitar directamente. El mecanismo es quitarle a la métrica un grado de libertad: en lugar de dejar que el elemento de volumen $\sqrt{-g}$ sea dinámico, se lo fija a una función dada.

Históricamente eso se hizo restringiendo la teoría a métricas con $\sqrt{-g}=1$ en algún sistema de coordenadas (y, en consecuencia, a transformaciones de coordenadas de jacobiano unidad); las referencias clásicas son \cite{anderson71,ng91}, que no se leyeron para este trabajo. Una formulación completamente covariante, debida a Henneaux y Teitelboim \cite{henneaux89}, reemplaza la condición $\sqrt{-g}=1$ por un campo auxiliar adicional; su acción, según la reproduce \cite{padilla15}, es una de las formas equivalentes de la teoría. Aquí se usa la variante con \emph{multiplicador de Lagrange} y una densidad de volumen fiducial, que es la de \cite{bengochea23}, porque hace explícito en las ecuaciones dónde entra el vínculo.

\subsection{La acción con multiplicador de Lagrange}

Se toma la acción de Einstein--Hilbert y se agrega un término que impone que el volumen de la métrica coincida con un volumen fijo dado \etq{C},
\begin{equation}
S=\frac{1}{2\kap}\int\dd^4x\,\Big[\sqrt{-g}\,R-2\lambda(x)\big(\sqrt{-g}-f(x)\big)\Big]+S_m[g_{ab},\Psi].
\label{eq:SUG}
\end{equation}
La función $f(x)$ es la densidad de una 4-forma de volumen fiducial $\varepsilon_{abcd}$: un dato de la teoría, que no es dinámico. El campo $\lambda(x)$ es un multiplicador de Lagrange. Como analogía (no forma parte de la derivación): en un péndulo con varilla rígida, la tensión de la varilla es el multiplicador del vínculo ``largo fijo''; nadie la elige, la determina el movimiento. Acá $\lambda$ cumple ese papel para el vínculo ``volumen fijo''.

Las ecuaciones de campo salen de variar respecto de $\lambda$, de la métrica y de la materia. Variar respecto de $\lambda$ da directamente el vínculo,
\begin{equation}
\sqrt{-g}=f.
\label{eq:vinculo}
\end{equation}
Para variar respecto de $g^{ab}$ se usa $\delta\sqrt{-g}=-\tfrac12\sqrt{-g}\,g_{ab}\,\delta g^{ab}$, de modo que el término del vínculo aporta $-2\lambda\,\delta\sqrt{-g}=+\lambda\sqrt{-g}\,g_{ab}\,\delta g^{ab}$, y con $\delta S_m=-\tfrac12\int\sqrt{-g}\,T_{ab}\,\delta g^{ab}$ queda $\frac{1}{2\kap}\sqrt{-g}\,[G_{ab}+\lambda g_{ab}-\kap T_{ab}]\,\delta g^{ab}=0$. Es decir \etq{C, P},
\begin{equation}
G_{ab}+\lambda(x)\,g_{ab}=\kap\,T_{ab}.
\label{eq:EinsteinUG}
\end{equation}
Las ecuaciones \eqref{eq:EinsteinUG} y \eqref{eq:vinculo} son el sistema de la teoría, junto con las ecuaciones de movimiento de la materia. La comparación con \eqref{eq:EinsteinLambda} es reveladora: la ecuación tiene la misma forma que la de RG con constante cosmológica, pero acá el coeficiente $\lambda$ \emph{no es una constante de la teoría sino un campo}, una incógnita más. Un detalle que confunde a veces: el vínculo \eqref{eq:vinculo} no dice ``$\sqrt{-g}=1$'', sino $\sqrt{-g}=f$, y solo se reduce a $\sqrt{-g}=1$ si se eligen coordenadas en las que $f=1$.

\subsection{La traza y la energía del vacío}

Hay que entender bien qué se pierde y qué se gana. Tomando la traza de \eqref{eq:EinsteinUG}, con $g^{ab}G_{ab}=-R$ y $g^{ab}g_{ab}=4$, se obtiene $-R+4\lambda=\kap T$, de donde
\begin{equation}
\lambda=\frac{R+\kap T}{4}.
\label{eq:lambdatraza}
\end{equation}
En RG con constante cosmológica, la traza es una ecuación dinámica para la geometría, $R=-\kap T+4\Lambda$. En UG la misma operación \emph{no da una ecuación para la métrica}: solo dice cuánto vale $\lambda$ una vez conocidas $R$ y $T$. Al reemplazar $\lambda$ en \eqref{eq:EinsteinUG} queda únicamente la parte sin traza,
\begin{equation}
R_{ab}-\tfrac14\,R\,g_{ab}=\kap\Big(T_{ab}-\tfrac14\,T\,g_{ab}\Big).
\label{eq:trazafree}
\end{equation}
Son las llamadas ecuaciones de Einstein sin traza \cite{ellis11,ellis14}: nueve ecuaciones independientes en lugar de diez. Esa es la ecuación que se perdió, y el vínculo \eqref{eq:vinculo} es lo que ocupa su lugar.

La ganancia se ve en \eqref{eq:trazafree}. Si en $T_{ab}$ se suma un término $c\,g_{ab}$ con $c$ constante (que es lo que hace una energía de vacío), la traza cambia en $4c$ y el lado derecho de \eqref{eq:trazafree} no cambia: $(T_{ab}+cg_{ab})-\tfrac14(T+4c)g_{ab}=T_{ab}-\tfrac14Tg_{ab}$ \etq{P}. En términos de \eqref{eq:EinsteinUG}, el efecto es solo desplazar $\lambda\to\lambda+\kap c$. La energía del vacío \emph{no gravita directamente}; queda absorbida en el multiplicador. Esta es la motivación de la teoría frente al problema de la constante cosmológica. Es importante no sobrevalorarla. UG no predice el valor de la constante cosmológica observada: como se verá en la sección siguiente, sigue apareciendo una constante de integración $\Lambda_0$ cuyo valor no está fijado por la teoría, y \cite{ellis11} discute justamente que las ecuaciones sin traza no resuelven el problema del valor de $\Lambda$ (lo que solo se leyó en el resumen del artículo).

\subsection{La conservación deja de ser un teorema: la difusión $Q$}
\label{sec:difusionQ}

Aquí está el punto de la teoría que da lugar a todo lo que sigue. Se aplica $\nabla^a$ a \eqref{eq:EinsteinUG} y se usa la identidad de Bianchi, $\nabla^aG_{ab}=0$, \emph{sin suponer nada} sobre la conservación de $T_{ab}$. Como $\nabla^a(\lambda g_{ab})=\nabla_b\lambda$, resulta
\begin{equation}
\nabla_b\lambda=\kap\,\nabla^aT_{ab}.
\label{eq:gradlambda}
\end{equation}
Esta ecuación cambia por completo el estatus de la conservación. En RG era una consecuencia automática de las ecuaciones de campo, \eqref{eq:conservacion}. En UG, la conservación de $T_{ab}$ equivale a que $\lambda$ sea constante, y se convierte en una hipótesis adicional sobre la materia. Hay dos casos.

Si la materia se conserva, $\nabla^aT_{ab}=0$, entonces $\nabla_b\lambda=0$ y $\lambda=\Lambda_0$ es una constante. La ecuación \eqref{eq:EinsteinUG} coincide con \eqref{eq:EinsteinLambda}: es la RG con constante cosmológica $\Lambda=\Lambda_0$, pero con una diferencia conceptual decisiva. Esa constante \emph{no está en la acción}, es una constante de integración: distintos valores de $\Lambda_0$ no son distintas teorías sino distintas soluciones (distintos datos iniciales) de la misma teoría. Esta equivalencia clásica con la RG en el caso conservativo es un resultado bien establecido de la literatura \cite{padilla15,ellis11,bengochea23,carballo22}.

Si la materia no se conserva, $\nabla^aT_{ab}\neq0$, la ecuación \eqref{eq:gradlambda} dice que el vector $\nabla^aT_{ab}$ es un gradiente (el de $\lambda/\kap$). Es natural definir un campo escalar $Q(x)$ por la relación
\begin{equation}
\nabla_aT^{ab}=\nabla^bQ,\qquad \lambda=\Lambda_0+\kap\,Q(x),
\label{eq:Qdef}
\end{equation}
donde $\Lambda_0$ es una constante de integración \etq{C: \cite{bengochea23,leon22}}. En lo que sigue, $Q$ se llama \emph{difusión}. La primera ecuación de \eqref{eq:Qdef} equivale a que el tensor $T^{ab}-g^{ab}Q$ se conserve, $\nabla_a(T^{ab}-g^{ab}Q)=0$. Eso permite una interpretación física útil (que se verifica en la sección~\ref{sec:fondo}): $Q$ actúa sobre la geometría como un fluido adicional con densidad $\rho_Q=Q$ y presión $p_Q=-Q$, o sea con ecuación de estado $w=-1$, pero con una densidad que \emph{varía} con el tiempo, y que intercambia energía con la materia ordinaria. El campo $Q$ es un escalar por construcción, y esto será lo que impida que afecte a las perturbaciones tensoriales.

\begin{alerta}[Qué significa que la materia no se conserve]
Conviene ser explícito sobre lo que implica \eqref{eq:Qdef}. Si la materia se describe con una acción invariante frente a difeomorfismos y satisface sus ecuaciones de movimiento, entonces $\nabla^aT_{ab}=0$ se cumple de manera idéntica y $Q$ es una constante. Un $Q$ que varía exige, por lo tanto, que la dinámica de la materia \emph{no sea} la de una acción covariante estándar: por ejemplo, por los efectos de gravedad cuántica o de modificaciones no unitarias de la mecánica cuántica que se invocan en \cite{josset17,perez19}. En los modelos de inflación de \cite{leon22,piccirilli23} esto se implementa de manera fenomenológica: se postula la ecuación de no conservación y se estudian sus consecuencias. Es una hipótesis sobre la materia, no un resultado de la teoría, y condiciona la lectura de todo lo que sigue (véase especialmente la sección~\ref{sec:supuestos}).
\end{alerta}

El origen físico propuesto para $Q$ merece una palabra. Josset, Perez y Sudarsky \cite{josset17} argumentan que en UG la violación de la conservación de la energía, debida por ejemplo a modificaciones no unitarias de la mecánica cuántica o a modelos de discretización del espacio-tiempo, genera un efecto equivalente a una constante cosmológica efectiva; Perez y Sudarsky \cite{perez19} desarrollan la idea desde la discreción de la gravedad cuántica. Ninguno de los dos artículos se leyó en detalle (solo el resumen del primero y los metadatos del segundo), de modo que lo anterior recoge solo lo que dicen sus resúmenes. Landau y colaboradores \cite{landau23} confrontan el fondo cosmológico con difusión con datos, y un comentario crítico \cite{cataldo26} sostiene, según su resumen, que la segunda ley de la termodinámica exige $\dot Q<0$. Esa condición se respeta en el modelo numérico de este trabajo, donde $Q$ decrece.

\subsection{Coordenadas, tiempo unimodular y simetrías}

Como \eqref{eq:EinsteinUG} y \eqref{eq:vinculo} son ecuaciones covariantes, el vínculo no restringe las soluciones: si se conoce una solución en unas coordenadas, basta expresar en ellas la densidad de volumen fiducial para que la solución satisfaga \eqref{eq:vinculo}. Para el universo homogéneo, en tiempo cósmico ($\dd s^2=-\dd t^2+a^2\dd\mathbf x^2$) es $\sqrt{-g}=a^3$, de modo que $f=a^3$. La coordenada en la que se cumple $\sqrt{-g}=1$ es el \emph{tiempo unimodular} $T$, definido por $\dd T=a^3\dd t$; en ella la misma solución se escribe
\begin{equation}
\dd s^2=-a^{-6}\,\dd T^2+a^2\,\dd\mathbf x^2.
\label{eq:metricaT}
\end{equation}
Es la misma física reetiquetada \cite{bengochea23}. Trabajar en tiempo cósmico, como se hace en este informe, equivale a una elección de coordenadas y no cambia ningún resultado (sección~\ref{sec:tensores} incluye la ecuación de los modos en $T$ para mostrarlo).

Sí hay una consecuencia de la estructura de UG que afecta a las perturbaciones: al haber fijado el volumen, los cambios de coordenadas admisibles (los que preservan la forma del vínculo) están restringidos a los de divergencia nula, $\nabla_a\xi^a=0$. Ese es el origen de las discrepancias que se discuten a continuación sobre el gauge de las perturbaciones escalares.

\subsection{Qué dice la literatura sobre perturbaciones y ondas gravitacionales}
\label{sec:ug-literatura}

Se resume lo que se pudo leer, con la advertencia de que de unos treinta trabajos relevantes solo se leyó el texto (completo o las secciones clave) de unos diez, y del resto solo el resumen (Apéndice~\ref{ap:fuentes}).

Para el caso conservativo hay consenso: UG equivale a RG con la constante cosmológica como constante de integración \cite{padilla15,ellis11,bengochea23,carballo22}. A nivel lineal, los autores de \cite{carballo22} encuentran que UG corresponde a una teoría llamada WTDiff y RG a la de Fierz--Pauli, sin otras diferencias que el tratamiento de la constante cosmológica, según el resumen de su revisión. Para las perturbaciones cosmológicas, Basak, Fabre y Shankaranarayanan \cite{basak16} afirman que las ecuaciones de fondo y de perturbación de UG y RG son idénticas, y que en particular, ``por ser sin traza y sin divergencia, los modos vectoriales y tensoriales no son afectados por el vínculo unimodular'', de modo que su evolución es idéntica. Esta afirmación, sin embargo, es para materia conservada. Gao, Brandenberger, Cai y Chen \cite{gao14} sostienen que el vínculo restringe el gauge escalar y que eso tiene consecuencias (por ejemplo, para el efecto Sachs--Wolfe), y trabajos posteriores \cite{alvarenga23,linares23} encuentran diferencias en el sector escalar; Basak et al. atribuyen esas diferencias a la elección de gauge. Nadie de los trabajos leídos reporta diferencias en el sector tensorial.

Sobre los tensores en concreto hay cuatro resultados relevantes. Cho y Singh \cite{cho15} estudian una UG \emph{generalizada} con un parámetro $\xi$ (la RG corresponde a $\xi=6$) y muestran que la acción tensorial no depende de $\xi$: ``debido a la ausencia de traza de la perturbación tensorial, $\sqrt{-g}$ queda fijo hasta primer orden, y no afecta a la ecuación del modo tensorial''; su espectro $P_T$ es el estándar, y $r$ cambia respecto de RG solo a través del fondo. Fabris y colaboradores \cite{fabris22} es el único trabajo encontrado que estudia ondas gravitacionales en UG \emph{no conservativa}: obtienen una ecuación que es ``exactamente la de las ondas gravitacionales en RG con un fluido perfecto'' y cuyas diferencias con $\Lambda$CDM provienen solo de las funciones de fondo $H$ y $\dot H$, aunque advierten que su análisis es esencialmente cualitativo y que el modelo necesita una condición adicional para cerrar las ecuaciones. Chakraborty, Jana y Mohanty \cite{chakraborty25} estudian la emisión de ondas por sistemas binarios en UG y encuentran que la suma sobre polarizaciones de los gravitones sin masa es la de RG. Y Barvinsky y Kolganov \cite{barvinsky19} muestran que en una UG generalizada, que agrega un gravitón escalar, la acción de los tensores tiene la forma de la de RG y que $r$ cambia solo por el sector escalar.

\begin{resultado}[El hueco que ocupa este trabajo]
Los resultados anteriores apuntan todos en la misma dirección, pero ninguno de los que se leyeron deriva, partiendo de la acción con multiplicador de Lagrange y con difusión $Q$, la ecuación de los modos tensoriales y la acción cuadrática \emph{incluyendo el término del vínculo a segundo orden}, y ninguno calcula $P_T(k)$ para un inflatón con difusión. Además, los trabajos que sí tratan inflación con difusión \cite{leon22,piccirilli23} asumen que el sector tensorial es el de RG: \cite{leon22} no contiene sector tensorial, y \cite{piccirilli23} da $r=16\epsilon_1$ sin derivarlo. Ese es el hueco que este trabajo intenta llenar en su parte tensorial.
\end{resultado}
